% GENERATED FILE. Edit canonical lessons, then run npm run generate:books.
% canonical-source-hash: fnv1a64:299c8cf6eb16ddac
% canonical-lessons: GU-C456-dubhavu, GU-C456-mathu, GU-C456-mahenu, GU-C456-avganna, GU-C456-thes

\chapter{Hurt Feelings}
\label{ch:gu-hurt-feelings}
\hlchaptermodality{456}

\begin{chapteropening}
\textbf{By the end of this chapter:} \emph{I can say when words hurt someone.}

Complete the last lesson of chapter 456: I can say when words hurt someone.
\end{chapteropening}

\section[\texorpdfstring{\gu{દુભાવું} (dubhāvũ)}{દુભાવું (dubhāvũ)}]{\gu{દુભાવું} (dubhāvũ) — to be hurt (in one's feelings)}

\label{lesson:GU-C456-dubhavu}



\begin{quote}
Before the new one: say the Gujarati for to make someone cry, then the Gujarati for weeping.
\end{quote}

\subsection*{You'll want to know: \gu{દુભાવું}}
\textbf{\gu{દુભાવું}} — \emph{dubhāvũ} — \textquotedblleft{}to be hurt (in one's feelings)\textquotedblright{}.

\begin{quote}
\gu{મારા} \gu{જવાબથી} \gu{મીરા} \gu{દુભાઈ} \gu{ગઈ}.
\end{quote}

\begin{quote}
\emph{mārā javābthī Mīrā dubhāī gaī.}
\end{quote}

\begin{quote}
Mira was hurt by my answer.
\end{quote}

Past (she …): \gu{તે} \gu{દુભાઈ} \gu{ગઈ}, \emph{te dubhāī gaī}.

\begin{cousinweb}
\textbf{\gu{દુભાવું}} (\emph{dubhāvũ}) is inherited from Sanskrit \emph{dabhnōti}, \textquotedblleft{}hurts, deceives\textquotedblright{}; CDIAL lists the Gujarati verbs of this family, \textquotedblleft{}to tease, vex\textquotedblright{} (CDIAL 6175).
\end{cousinweb}

\begin{grammarlens}[title={What you've built}]
One more word for describing what you lived and felt.
\end{grammarlens}

\subsection*{Your turn}
\begin{itemize}
\raggedright
  \item \emph{Say it:} \emph{dubhāvũ}
  \item \emph{Say it:} \emph{dubhāvũ}, once more
  \item \emph{Say it:} \emph{mārā javābthī Mīrā dubhāī gaī.} — Mira was hurt by my answer.
  \item \emph{Recall:} say the Gujarati for to make someone cry, then the Gujarati for weeping, then say \emph{dubhāvũ} again
\end{itemize}

\begin{tcolorbox}[breakable,colback=teal!4,colframe=teal!35!black,title={Before you move on}]
What does \gu{દુભાવું} mean? (\textquotedblleft{}To be hurt (in one's feelings)\textquotedblright{}.) Say it once more.
\end{tcolorbox}
\section[\texorpdfstring{\gu{માઠું} (māṭhũ)}{માઠું (māṭhũ)}]{\gu{માઠું} (māṭhũ) — unwelcome; bad (of news or feelings)}

\label{lesson:GU-C456-mathu}



\begin{quote}
Before the new one: say the Gujarati for weeping, then the Gujarati for to be hurt.
\end{quote}

\subsection*{You'll want to know: \gu{માઠું}}
\textbf{\gu{માઠું}} — \emph{māṭhũ} — \textquotedblleft{}unwelcome; bad (of news or feelings)\textquotedblright{}.

\begin{quote}
\gu{મારા} \gu{જવાબનું} \gu{તેને} \gu{માઠું} \gu{લાગ્યું}.
\end{quote}

\begin{quote}
\emph{mārā javābnũ tene māṭhũ lāgyũ.}
\end{quote}

\begin{quote}
He took offence at my answer.
\end{quote}

\begin{cousinweb}
\textbf{\gu{માઠું}} (\emph{māṭhũ}) goes back to a reconstructed \emph{*maṭṭa-}, \textquotedblleft{}defective\textquotedblright{}; CDIAL lists the Gujarati word, \textquotedblleft{}vicious, bad, unfortunate\textquotedblright{} (CDIAL 9723).
\end{cousinweb}

\begin{grammarlens}[title={What you've built}]
One more word for describing what you lived and felt.
\end{grammarlens}

\subsection*{Your turn}
\begin{itemize}
\raggedright
  \item \emph{Say it:} \emph{māṭhũ}
  \item \emph{Say it:} \emph{māṭhũ}, once more
  \item \emph{Say it:} \emph{mārā javābnũ tene māṭhũ lāgyũ.} — He took offence at my answer.
  \item \emph{Recall:} say the Gujarati for weeping, then the Gujarati for to be hurt, then say \emph{māṭhũ} again
\end{itemize}

\begin{tcolorbox}[breakable,colback=teal!4,colframe=teal!35!black,title={Before you move on}]
What does \gu{માઠું} mean? (\textquotedblleft{}Unwelcome; bad (of news or feelings)\textquotedblright{}.) Say it once more.
\end{tcolorbox}
\section[\texorpdfstring{\gu{મહેણું} (maheṇũ)}{મહેણું (maheṇũ)}]{\gu{મહેણું} (maheṇũ) — a taunt}

\label{lesson:GU-C456-mahenu}



\begin{quote}
Before the new one: say the Gujarati for to be hurt, then the Gujarati for unwelcome.
\end{quote}

\subsection*{You'll want to know: \gu{મહેણું}}
\textbf{\gu{મહેણું}} — \emph{maheṇũ} — \textquotedblleft{}a taunt\textquotedblright{}.

\begin{quote}
\gu{પડોશીએ} \gu{તેને} \gu{મહેણું} \gu{માર્યું}.
\end{quote}

\begin{quote}
\emph{paḍośīe tene maheṇũ māryũ.}
\end{quote}

\begin{quote}
The neighbour taunted her.
\end{quote}

A neuter noun.

\begin{cousinweb}
\textbf{\gu{મહેણું}} (\emph{maheṇũ}) goes back to a reconstructed \emph{*mahākathana-}, \textquotedblleft{}big words\textquotedblright{}; CDIAL lists the Gujarati word, \textquotedblleft{}taunt, bitter remark\textquotedblright{} (CDIAL 9938).
\end{cousinweb}

\begin{grammarlens}[title={What you've built}]
One more word for describing what you lived and felt.
\end{grammarlens}

\subsection*{Your turn}
\begin{itemize}
\raggedright
  \item \emph{Say it:} \emph{maheṇũ}
  \item \emph{Say it:} \emph{maheṇũ}, once more
  \item \emph{Say it:} \emph{paḍośīe tene maheṇũ māryũ.} — The neighbour taunted her.
  \item \emph{Recall:} say the Gujarati for to be hurt, then the Gujarati for unwelcome, then say \emph{maheṇũ} again
\end{itemize}

\begin{tcolorbox}[breakable,colback=teal!4,colframe=teal!35!black,title={Before you move on}]
What does \gu{મહેણું} mean? (\textquotedblleft{}A taunt\textquotedblright{}.) Say it once more.
\end{tcolorbox}
\section[\texorpdfstring{\gu{અવગણના} (avgaṇnā)}{અવગણના (avgaṇnā)}]{\gu{અવગણના} (avgaṇnā) — disregard; being ignored}

\label{lesson:GU-C456-avganna}



\begin{quote}
Before the new one: say the Gujarati for unwelcome, then the Gujarati for a taunt.
\end{quote}

\subsection*{You'll want to know: \gu{અવગણના}}
\textbf{\gu{અવગણના}} — \emph{avgaṇnā} — \textquotedblleft{}disregard; being ignored\textquotedblright{}.

\begin{quote}
\gu{માતા} \gu{અને} \gu{પિતાની} \gu{અવગણના} \gu{ન} \gu{કરવી} \gu{જોઈએ}.
\end{quote}

\begin{quote}
\emph{mātā ane pitānī avgaṇnā na karvī joīe.}
\end{quote}

\begin{quote}
One should not neglect one's mother and father.
\end{quote}

A feminine noun.

\begin{cousinweb}
\textbf{\gu{અવગણના}} (\emph{avgaṇnā}) is a tatsama, close to Sanskrit \emph{avagaṇana}, \textquotedblleft{}disregard, contempt\textquotedblright{}, as Monier-Williams records it, with the feminine \emph{-ā}.
\end{cousinweb}

\begin{grammarlens}[title={What you've built}]
One more word for describing what you lived and felt.
\end{grammarlens}

\subsection*{Your turn}
\begin{itemize}
\raggedright
  \item \emph{Say it:} \emph{avgaṇnā}
  \item \emph{Say it:} \emph{avgaṇnā}, once more
  \item \emph{Say it:} \emph{mātā ane pitānī avgaṇnā na karvī joīe.} — One should not neglect one's mother and father.
  \item \emph{Recall:} say the Gujarati for unwelcome, then the Gujarati for a taunt, then say \emph{avgaṇnā} again
\end{itemize}

\begin{tcolorbox}[breakable,colback=teal!4,colframe=teal!35!black,title={Before you move on}]
What does \gu{અવગણના} mean? (\textquotedblleft{}Disregard; being ignored\textquotedblright{}.) Say it once more.
\end{tcolorbox}
\section[\texorpdfstring{\gu{ઠેસ} (ṭhes)}{ઠેસ (ṭhes)}]{\gu{ઠેસ} (ṭhes) — a knock; a hurt to one's pride}

\label{lesson:GU-C456-thes}



\begin{quote}
Before the new one: say the Gujarati for a taunt, then the Gujarati for disregard.
\end{quote}

\subsection*{You'll want to know: \gu{ઠેસ}}
\textbf{\gu{ઠેસ}} — \emph{ṭhes} — \textquotedblleft{}a knock; a hurt to one's pride\textquotedblright{}.

\begin{quote}
\gu{તેના} \gu{જવાબથી} \gu{મને} \gu{ઘણી} \gu{ઠેસ} \gu{લાગી}.
\end{quote}

\begin{quote}
\emph{tenā javābthī mane ghaṇī ṭhes lāgī.}
\end{quote}

\begin{quote}
His answer was a hard knock for me.
\end{quote}

A feminine noun.

\begin{cousinweb}
\textbf{\gu{ઠેસ}} (\emph{ṭhes}) goes back to a reconstructed \emph{*ṭhēcc-}, \textquotedblleft{}push\textquotedblright{}; CDIAL lists the Gujarati word, \textquotedblleft{}knock\textquotedblright{} (CDIAL 5511), and Gujarati also uses it for a knock to the feelings.
\end{cousinweb}

\begin{grammarlens}[title={What you've built}]
One more word for describing what you lived and felt.
\end{grammarlens}

\subsection*{Your turn}
\begin{itemize}
\raggedright
  \item \emph{Say it:} \emph{ṭhes}
  \item \emph{Say it:} \emph{ṭhes}, once more
  \item \emph{Say it:} \emph{tenā javābthī mane ghaṇī ṭhes lāgī.} — His answer was a hard knock for me.
  \item \emph{Recall:} say the Gujarati for a taunt, then the Gujarati for disregard, then say \emph{ṭhes} again
\end{itemize}

\begin{tcolorbox}[breakable,colback=teal!4,colframe=teal!35!black,title={Before you move on}]
What does \gu{ઠેસ} mean? (\textquotedblleft{}A knock; a hurt to one's pride\textquotedblright{}.) Say it once more.
\end{tcolorbox}
